\chapter{Граница углового разложения центрированного вихря}
\label{ch:v6-bosonic-defect-centered-angular-channel-decomposition-gate}

\section{Проблема}

После вынесения двух переносов в коллективные координаты остаётся кластер
малых положительных уровней. Чтобы решить, являются ли они внутренними модами
или ещё одним следствием дискретизации, требуется разложить возмущения по
угловым каналам. Для осесимметричного непрерывного вихря такой переход
каноничен: линеаризованный оператор коммутирует с вращениями и распадается на
независимые радиальные блоки \cite{alonso-vortex-spectrum,
ginzburg-landau-morse-index}.

\section{Проверяемая схема}

На сетке $33\times33$ заново построено центрированное стационарное решение.
Среднеквадратичная норма проецированного градиента равна
\[
  4.20\cdot10^{-9}.
\]
Вычислены первые двадцать четыре собственных значения. Они лежат между
$0.0008661276$ и $0.0113575006$; отрицательных уровней в этом окне нет.

Далее построен точный оператор четверть-поворота $R$. Он одновременно
вращает координаты, комплексную заряженную амплитуду и две компоненты
связности. Алгебраическое тождество
\[
  R^4=1
\]
выполняется с остатком $1.2\cdot10^{-16}$.

\section{Результат}

Из восьми согласованных комбинаций направления пространственного поворота и
знаков внутренних действий выбрана та, которая лучше всего сохраняет
стационарный фон. Даже для неё относительный остаток равен
\[
  \frac{\|R\Phi_0-\Phi_0\|}{\|\Phi_0\|}=0.0327391.
\]
Первые двадцать четыре моды не образуют представление поворотов:
средняя утечка из их подпространства равна $0.679717$, максимальная ---
$0.928113$, а минимальное сингулярное число сжатого оператора поворота равно
$1.47\cdot10^{-7}$. Максимальный относительный остаток коммутатора гессиана с
поворотом равен $0.999972$.

\begin{proposition}
Текущая односторонняя декартова дискретизация не сертифицирует угловые метки
непрерывного вихря. Отсутствие отрицательных значений среди двадцати четырёх
нижних уровней сохраняется как конечносеточный факт, но распределять эти
уровни по угловым моментам на основании данного оператора нельзя.
\end{proposition}

Механизм нарушения прозрачен. При повороте на $\pi/2$ односторонняя разность
в одном направлении переходит в разность противоположной ориентации. Такая
схема была выбрана ранее, чтобы не вводить решёточные удвоители центральной
разности; цена выбора --- потеря точной вращательной ковариантности.

\section{Нулевая гипотеза и критерий остановки}

Нулевая гипотеза: после перехода к полярной постановке мягкий кластер
разделится на переносный канал $|m|=1$ и неотрицательные внутренние блоки.
Первичная литература подтверждает, что переносы лежат именно в первом
угловом канале, а возможные внутренние и распадные моды классифицируются
отдельными мультипольными блоками \cite{sigal-vortex-decomposition,
alonso-vortex-spectrum}.

Стоп-критерий следующего гейта: сходящееся отрицательное собственное значение
любого полярного блока закрывает прямую вихревую нить. Если отрицательных
значений нет, но нижняя граница снова стремится к нулю вне переносного блока,
положительный внутренний зазор также не считается доказанным.

\section{Статус}

Текущий декартов расчёт не обнаружил отрицательной моды, но угловая
классификация на нём отвергнута. Следующий обязательный шаг --- прямой
полярный оператор Штурма--Лиувилля для каждого углового числа. Полный гессиан
$Q+T+B$, замыкание нити и рождение наблюдаемой материи остаются открытыми.

\begin{thebibliography}{9}\raggedright
\bibitem{alonso-vortex-spectrum}
A.~Alonso-Izquierdo, D.~Migu\'elez-Caballero, J.~Queiruga,
\emph{Spectral structure of fluctuations around $n$-vortices in the
Abelian-Higgs model}, arXiv:2505.05039.

\bibitem{ginzburg-landau-morse-index}
\emph{Nondegeneracy and Morse Index of Ginzburg--Landau Vortices},
arXiv:2608.04976.

\bibitem{sigal-vortex-decomposition}
I.~M.~Sigal,
\emph{Magnetic Vortices, Abrikosov Lattices and Automorphic Functions},
arXiv:1308.5440.
\end{thebibliography}